% ============================================================
% Dark Store Placement + Integrated Forward‑Reverse Logistics
% Project Report — Day 7 Draft
% Generated by src/report_builder.py
% ============================================================
\documentclass[11pt,a4paper]{article}

\usepackage[utf8]{inputenc}
\usepackage[T1]{fontenc}
\usepackage{lmodern}
\usepackage[margin=2.4cm]{geometry}
\usepackage{booktabs}
\usepackage{graphicx}
\usepackage{float}
\usepackage{amsmath}
\usepackage{amssymb}
\usepackage{hyperref}
\usepackage{setspace}
\usepackage{enumitem}
\usepackage{caption}
\usepackage{subcaption}
\usepackage{xcolor}
\usepackage{parskip}
\usepackage{microtype}

\definecolor{darkblue}{HTML}{1565C0}
\hypersetup{colorlinks=true, linkcolor=darkblue, urlcolor=darkblue, citecolor=darkblue}

\setstretch{1.15}

% ── Title ──────────────────────────────────────────────────────────────────
\title{
  \textbf{Dark Store Placement \&\\
  Integrated Forward--Reverse Logistics Optimisation}\\[0.4em]
  \large Project Report --- PGDBA / ISI Kolkata / IIM Calcutta / IIT Kharagpur\\[0.2em]
  \normalsize Operations Management $\cdot$ Supply Chain Analytics
}
\author{
  Pritam (Lead Coder \& Integration Architect) \\
  Sneha $\cdot$ Vybhav $\cdot$ Pranav $\cdot$ Team\\[0.2em]
  \small Dataset: Olist Brazilian E-Commerce (Kaggle)
}
\date{April 2026}

\begin{document}
\maketitle
\tableofcontents
\newpage

% ===========================================================================
\section{Introduction}
% ===========================================================================

E-commerce last-mile logistics in dense urban markets faces two coupled
challenges typically treated independently: \textbf{forward delivery} of
orders from dark stores to customers, and \textbf{reverse logistics} ---
collecting returned items without deploying a separate fleet.

This project asks a single integrated question:
\begin{quote}
  \textit{Given customer locations, order demands, and predicted return
  probabilities, (a) where should dark stores be placed to minimise
  customer-to-facility distance, and (b) how should vehicle routes be designed
  to simultaneously deliver orders and collect returns, minimising a weighted
  combination of total cost, delivery latency, and fleet size?}
\end{quote}

The solution pipeline consists of five coupled stages implemented over eight
days using Google OR-Tools~9.15, PuLP~3.3 (CBC solver), and XGBoost.
All production logic resides in \texttt{src/}; notebooks are used only for
testing and visualisation.

% ===========================================================================
\section{Dataset and Environment}
% ===========================================================================

\subsection{Olist Brazilian E-Commerce}
The Olist dataset contains $\approx$100,000 orders across 9 CSV tables
(orders, items, products, customers, sellers, geolocation, payments,
reviews, category translation), covering 2016--2018.
We filter to \textbf{customer\_state = `SP'} (S\~ao Paulo, $\approx$41\% of
all orders), yielding \textbf{19,207 rows} for the active pipeline
after preprocessing, geolocation resolution, and feature engineering.

\subsection{Environment}
\begin{itemize}[noitemsep]
  \item Python 3.13.12 managed by \texttt{uv} on WSL2
  \item OR-Tools 9.15.6755 --- CVRPTW + SDVRP routing
  \item PuLP 3.3.0 + CBC --- MILP joint optimiser
  \item XGBoost 3.2.0 + Platt calibration --- return classifier
  \item All code: \url{https://github.com/metaphorpritam/SCA_DARK_STORES}
\end{itemize}

% ===========================================================================
\section{Stage 1: Dark Store Placement (Facility Location)}
% ===========================================================================

\subsection{K-Means with Demand Weighting}
Customer coordinates are weighted by zip-code order volume and clustered
using K-Means for $K \in [3, 12]$.  The optimal $K$ is selected at the
agreement point between the elbow curve (inertia) and silhouette score.

\begin{figure}[H]
  \centering
  \includegraphics[width=0.9\linewidth]{/mnt/d/Python-UV/SCA_DARK_STORES/outputs/elbow_silhouette.png}
  \caption{K-Means elbow + silhouette: optimal $K=11$.}
  \label{fig:elbow}
\end{figure}

\subsection{Coverage Rule and Final $K$}
Although the silhouette score peaks at $K=3$, the coverage rule
(70\% of customers within 5~km of their assigned dark store) requires
$K \geq 10$.  \textbf{$K = 11$} is selected, achieving 73.7\% coverage
within 5~km across S\~ao Paulo.

\begin{figure}[H]
  \centering
  \includegraphics[width=0.9\linewidth]{/mnt/d/Python-UV/SCA_DARK_STORES/outputs/kmeans_cluster_map.png}
  \caption{Final 11 dark store locations (K-Means centroids).}
  \label{fig:clusters}
\end{figure}

\subsection{p-Median Validation}
A PuLP MILP p-median formulation validates K-Means placement.
The objective minimises $\sum_i \sum_j d_{ij} \cdot w_i \cdot x_{ij}$
subject to single-assignment and open-facility constraints.
Both methods agree on cluster locations within 1~km, confirming stability.

% ===========================================================================
\section{Stage 2: Forward VRP --- CVRPTW (Day 3)}
% ===========================================================================

\subsection{Problem Formulation}
Each of the 11 zones is solved independently as a
\textbf{Capacitated VRP with Time Windows (CVRPTW)}:
\begin{align}
  \min &\sum_k \sum_{(i,j) \in A} c_{ij} \cdot x_{ijk} \notag\\
  \text{s.t.} \quad & \sum_k x_{0jk} = K, \quad
    \sum_j x_{ijk} = \sum_j x_{jik} \;\; \forall i,k \notag\\
  & a_i \leq s_{ik} \leq b_i, \quad
    \sum_i d_i \cdot \delta_{ik} \leq Q \;\; \forall k \notag
\end{align}
Parameters: vehicle capacity $Q = 500{\,}000\,\text{g}$,
speed $= 40\,\text{km/h}$, service time $= 5\,\text{min}$,
fixed cost R\$50/route, variable R\$1.50/km.
Solver: \texttt{PATH\_CHEAPEST\_ARC} $\to$ \texttt{GUIDED\_LOCAL\_SEARCH},
30~s time limit.

\subsection{Results --- All 11 Zones Solved}

\begin{tabular}{rrrrc}
      \toprule
      \textbf{Zone} & \textbf{Orders} & \textbf{Vehicles} &
      \textbf{Dist (km)} & \textbf{Cost (R\$)} \\
      \midrule
        0 & 75 & 2 & 94.00 & R\$241.00 \\
1 & 75 & 2 & 81.54 & R\$222.31 \\
2 & 75 & 1 & 108.01 & R\$212.02 \\
3 & 75 & 1 & 99.37 & R\$199.06 \\
4 & 75 & 3 & 88.57 & R\$282.86 \\
5 & 75 & 2 & 112.00 & R\$268.00 \\
6 & 75 & 2 & 97.33 & R\$246.00 \\
7 & 75 & 2 & 99.25 & R\$248.88 \\
8 & 75 & 3 & 113.91 & R\$320.86 \\
9 & 75 & 2 & 92.39 & R\$238.58 \\
10 & 75 & 2 & 83.22 & R\$224.83 \\
      \midrule
      \textbf{Total} & \textbf{825} &
      \textbf{22} & \textbf{1069.59} &
      \textbf{R\$2704.40} \\
      \bottomrule
      \end{tabular}

\vspace{0.5em}
\begin{center}
  \textbf{Total: 825 orders, 22 vehicles, 1069.59\,km,
  R\$2704.40}
\end{center}

\begin{figure}[H]
  \centering
  \includegraphics[width=0.9\linewidth]{/mnt/d/Python-UV/SCA_DARK_STORES/outputs/day3_zone_kpi_bars.png}
  \caption{Forward route cost and distance by zone.}
  \label{fig:fwdkpi}
\end{figure}

\subsection{Baseline Comparison}
Naive nearest-store routing (each customer drives directly to closest store)
produces 75,066.4~km.  OR-Tools CVRPTW achieves \textbf{1,069.59~km},
a \textbf{98.58\% reduction}.

% ===========================================================================
\section{Stage 3: Return Probability Classifier (Day 4)}
% ===========================================================================

\subsection{Target and Features}
Binary target: \texttt{is\_return = 1} if order status
$\in$ \{canceled, unavailable\} or delivered $>7$ days late.
Feature set: product weight, freight value, seller state (label-encoded),
review score, order value, days late, payment type, product category.

\subsection{Model: XGBoost + Platt Calibration}
Class imbalance ($\approx$5\% return rate) handled via
\texttt{scale\_pos\_weight = 15}.

\begin{tabular}{lc}
\toprule
\textbf{Metric} & \textbf{Value} \\
\midrule
ROC-AUC   & 0.8969 \\
PR-AUC    & 0.4716 \\
Brier score & 0.0213 \\
Precision @ 0.30 & 0.5556 \\
Recall @ 0.30    & 0.354 \\
F1 @ 0.30        & 0.4324 \\
\bottomrule
\end{tabular}

ROC-AUC~0.8969 exceeds the project target of $\geq 0.70$, confirming
the classifier's ability to identify return-risk orders.

Customers with \texttt{return\_prob > 0.30} are flagged as pickup nodes
in the SDVRP; expected pickup weight = \texttt{return\_prob $\times$
product\_weight\_g}.
\textbf{593 orders} (3.1\% of 19,207) were flagged across all zones.

% ===========================================================================
\section{Stage 4: Reverse VRP (Day 4)}
% ===========================================================================

Same CVRPTW structure as forward VRP with vehicles departing the depot
\emph{empty} and collecting returns (load increases at each stop).

\begin{tabular}{rrrrc}
      \toprule
      \textbf{Zone} & \textbf{Pickups} & \textbf{Vehicles} &
      \textbf{Dist (km)} & \textbf{Cost (R\$)} \\
      \midrule
        0 & 51 & 1 & 82.95 & R\$174.42 \\
1 & 75 & 2 & 81.48 & R\$222.22 \\
2 & 57 & 2 & 107.03 & R\$260.54 \\
3 & 34 & 2 & 59.31 & R\$188.96 \\
4 & 46 & 1 & 89.99 & R\$184.98 \\
5 & 42 & 1 & 93.72 & R\$190.58 \\
6 & 45 & 1 & 80.79 & R\$171.18 \\
7 & 53 & 1 & 75.86 & R\$163.79 \\
8 & 45 & 2 & 99.99 & R\$249.98 \\
9 & 57 & 1 & 87.80 & R\$181.70 \\
10 & 71 & 1 & 87.44 & R\$181.16 \\
      \midrule
      \textbf{Total} & \textbf{576} &
      \textbf{15} & \textbf{946.36} &
      \textbf{R\$2169.51} \\
      \bottomrule
      \end{tabular}

\vspace{0.5em}
\begin{center}
  \textbf{Total: 576 pickups, 15 vehicles,
  946.36\,km, R\$2169.51}
\end{center}

Combined separate cost:
\textbf{R\$4,873.91} (forward R\$2704.40 + reverse
R\$2169.51).

% ===========================================================================
\section{Stage 5: Joint MILP Optimiser (Day 4)}
% ===========================================================================

\subsection{Objective Function}
\begin{equation}
  Z = \alpha \cdot C_{\text{fwd}} + \beta \cdot C_{\text{rev}}
    + \gamma \cdot T_{\text{pen}} + \delta \cdot N_{\text{veh}}
  \label{eq:Z}
\end{equation}
Binary activation variables $u_v \in \{0,1\}$ (forward vehicles) and
$w_v \in \{0,1\}$ (reverse vehicles) are solved by PuLP CBC.

\subsection{Result ($\alpha=\beta=\gamma=\delta=0.25$)}
\begin{tabular}{lc}
\toprule
\textbf{Component} & \textbf{Value} \\
\midrule
$Z$ (objective)    & 54.38 \\
Status             & Optimal \\
$C_{\text{fwd}}$ & 44.942 \\
$C_{\text{rev}}$ & 169.577 \\
$T_{\text{pen}}$ & 287.063 \quad (\textbf{dominates $\approx$54\% of Z}) \\
$N_{\text{veh}}$ & 3 \\
\bottomrule
\end{tabular}

The penalty term $T_{\text{pen}}$ dominates $Z$ at equal weights,
motivating the Day~5 SDVRP work to reduce reverse fleet redundancy
and the Day~6 Pareto sweep to find balanced weight combinations.

% ===========================================================================
\section{Stage 6: SDVRP Hybrid Routing (Day 5)}
% ===========================================================================

\subsection{Motivation}
Zone~8 carries the highest combined cost: R\$570.84 (R\$320.86 forward
+ R\$249.98 reverse).  Deploying 5 separate vehicles (3 fwd + 2 rev)
means those vehicles pass the same customer locations twice.

\subsection{SDVRP Load Model}
Simultaneous Delivery and Pickup VRP (Dethloff,~2001) merges all
delivery and pickup nodes into one OR-Tools solve using a
\textbf{single load dimension}:
\begin{align}
  L(t) &= L_{\text{start}} - \text{delivered}(t) + \text{picked}(t) \notag\\
  \text{transit}[i] &= p_i - d_i \quad \text{(net change per stop)} \notag\\
  0 &\leq L_{\text{cumul}}[i] \leq Q \quad \forall\, i \notag
\end{align}
\texttt{fix\_start\_cumul\_to\_zero = False} lets OR-Tools initialise
each vehicle load at total delivery weight.  The previous two-dimension
additive approach ($d_{\text{cumul}} + p_{\text{cumul}} \leq Q$) was
over-constraining and is \textbf{corrected here}.

\subsection{All-Zone SDVRP Runner}
\texttt{run\_all\_zones\_sdvrp()} loops all 11 zones, calling
\texttt{solve\_sdvrp\_hybrid()} per zone and writing
\texttt{outputs/hybrid\_routes.json} and
\texttt{outputs/hybrid\_kpi\_summary.csv}.
Expected fleet saving: \textbf{15--25\%} vs separate fleets.

% ===========================================================================
\section{Stage 7: Combined KPI Report \& Pareto Sweep (Day 6)}
% ===========================================================================

\subsection{Combined KPI Report}
\texttt{src/kpi\_reporter.py} merges all three KPI streams
(forward, reverse, hybrid SDVRP) into
\texttt{outputs/combined\_kpi\_report.csv}.

\begin{figure}[H]
  \centering
  \includegraphics[width=0.9\linewidth]{/mnt/d/Python-UV/SCA_DARK_STORES/outputs/combined_cost_by_zone.png}
  \caption{Stacked fwd+rev cost by zone.  Zone~8 is largest at R\$570.84.}
  \label{fig:combined_cost}
\end{figure}

\subsection{Zone Priority Ranking}
Zones are ranked by \texttt{saving\_R\$} (when SDVRP results exist)
or by \texttt{separate\_cost\_R\$} otherwise.  Zone~8 receives
\textbf{Rank~1} in both cases.

\begin{figure}[H]
  \centering
  \includegraphics[width=0.9\linewidth]{/mnt/d/Python-UV/SCA_DARK_STORES/outputs/sdvrp_priority_ranking.png}
  \caption{SDVRP zone priority ranking (highest saving opportunity first).}
  \label{fig:priority}
\end{figure}

\subsection{Pareto Sweep}
Weight combinations $(\alpha, \beta) \in \{0.1, 0.3, 0.5, 0.7, 0.9\}^2$
with $\gamma = \delta = (1-\alpha-\beta)/2$ (valid for $\alpha+\beta \leq 1$)
yield up to 15 weight settings.  Pareto-optimal solutions minimise both
$C_{\text{routing}} = C_{\text{fwd}} + C_{\text{rev}}$ and
$T_{\text{pen}}$ simultaneously.

\begin{equation}
  \text{dist\_to\_ideal} =
  \sqrt{c_{\text{norm}}^2 + t_{\text{norm}}^2}
\end{equation}
The \textbf{knee point} (minimum distance to ideal) is the recommended
operating configuration.

\begin{figure}[H]
  \centering
  \includegraphics[width=0.9\linewidth]{/mnt/d/Python-UV/SCA_DARK_STORES/outputs/pareto_tradeoff.png}
  \caption{Pareto tradeoff surface: routing cost vs.\ late-return penalty.  Red star = knee point.}
  \label{fig:pareto}
\end{figure}

% ===========================================================================
\section{Summary of Results}
% ===========================================================================

\begin{tabular}{llr}
\toprule
\textbf{Stage} & \textbf{Metric} & \textbf{Value} \\
\midrule
Dark store placement   & Zones ($K$)                 & 11 \\
                       & Coverage within 5\,km       & 73.7\% \\
Forward VRP            & Total distance              & 1069.59\,km \\
                       & Improvement vs.\ naive      & 98.58\% \\
                       & Total cost                  & R\$2704.40 \\
Return classifier      & ROC-AUC                     & 0.8969 \\
                       & F1 @ threshold 0.30         & 0.4324 \\
Reverse VRP            & Total distance              & 946.36\,km \\
                       & Total cost                  & R\$2169.51 \\
Joint MILP             & $Z$ (optimal)               & 54.38 \\
Combined baseline      & Total fwd + rev cost        & R\$4,873.91 \\
\bottomrule
\end{tabular}

% ===========================================================================
\section{Reproducibility}
% ===========================================================================

All code is reproducible via \texttt{run\_all.sh}:
\begin{verbatim}
uv run python src/data_pipeline.py
uv run python src/clustering.py
uv run python src/forward_vrp.py
uv run python src/reverse_vrp.py
uv run python src/return_classifier.py
uv run python src/joint_optimizer.py
uv run python src/kpi_reporter.py
uv run python src/report_builder.py
\end{verbatim}

Data files (\texttt{data/raw/} Olist CSVs) are excluded from the repository
per \texttt{.gitignore}.  Download from Kaggle:
\url{https://www.kaggle.com/datasets/olistbr/brazilian-ecommerce}.

\section*{References}

\begin{enumerate}[noitemsep]
  \item Dethloff, J.\ (2001). Vehicle routing and reverse logistics:
        the vehicle routing problem with simultaneous delivery and pick-up.
        \textit{OR Spectrum}, 23(1), 79--96.
  \item Dantzig, G.B.\ \& Ramser, J.H.\ (1959). The truck dispatching problem.
        \textit{Management Science}, 6(1), 80--91.
  \item Google OR-Tools documentation:
        \url{https://developers.google.com/optimization}.
  \item ReVelle, C.S.\ \& Swain, R.W.\ (1970). Central facilities location.
        \textit{Geographical Analysis}, 2(1), 30--42.
\end{enumerate}

\end{document}
